\chapter*{Abstract}
\addcontentsline{toc}{chapter}{Abstract}

Autonomous Artificial Intelligence (AI) agents powered by Large Language Models (LLMs) are increasingly entrusted with privileged tool-calling capabilities, including querying sensitive internal databases, generating executable code, accessing cloud storage, and dispatching external electronic communications. However, modern transformer architectures inherently conflate passive untrusted data with privileged control instructions within a unified attention context. When autonomous agents ingest untrusted content from external documents, emails, or web pages, they become susceptible to Indirect Prompt Injection (OWASP LLM01). In such attacks, the adversary overrides the model's core objective, forcing it to act as an unauthenticated Confused Deputy (OWASP LLM06). Existing security paradigms—such as prompt engineering guardrails, input/output regex filtering, and naive Human-in-the-Loop (HITL) authorization—fail systematically because language models cannot police themselves, and human operators are readily deceived by socially engineered, plausible rationales generated by compromised agents.

To resolve this fundamental vulnerability, this capstone project introduces \textbf{AEGIS} (\textbf{Autonomous Enforcement Gateway for Information Security}), a zero-trust, in-line Policy Enforcement Point (PEP) proxy situated between autonomous AI agent runtimes and downstream tool brokers. AEGIS enforces a mathematical security foundation grounded in a dual-lattice Information Flow Control (IFC) model: Dorothy Denning's 1976 Confidentiality Lattice (employing Least Upper Bound join operations) and Kenneth Biba's 1977 Integrity Lattice (employing Greatest Lower Bound meet operations). By tracking fine-grained causal provenance across execution Directed Acyclic Graphs (DAGs), AEGIS establishes and enforces \textbf{Theorem 1 (Provenance Dominance Invariant)}, which mathematically guarantees that data taint strictly dominates human approval: no human approval token can authorize an action whose execution lineage contains untrusted inputs or confidential leaks.

AEGIS implements a 9-step Policy Enforcement Pipeline incorporating timing-safe HMAC agent authentication, dual-tier threat detection (sub-2ms heuristic regex filtering combined with local quantized ONNX transformer inference under 25ms), high-entropy Data Loss Prevention (DLP) secret detection across 8 regex families, canonical SHA-256 payload parameter binding, and an atomic single-use state machine. Concurrency is achieved via an optimized SQLite Write-Ahead Logging (WAL) architecture with graceful busy-timeout pooling, eliminating database lock contention. 

Extensive empirical evaluations validate the robustness, accuracy, and efficiency of AEGIS:
\begin{enumerate}
    \item \textbf{Unit \& Regression Verification:} A 106-test pytest suite spanning 19 modular files executes with a 100\% pass rate in 2.59 seconds.
    \item \textbf{Scenario Verification (E01--E10):} Across 10 standard adversarial and operational scenarios mapping to the OWASP Top 10 for LLMs, AEGIS achieves 100\% enforcement accuracy with a mean gateway latency of 17.98~ms (median 12.93~ms, min 5.65~ms, max 45.66~ms).
    \item \textbf{Repeated Reliability Benchmarks:} Across 300 repeated test executions (30 iterations per scenario), AEGIS demonstrates a 100\% decision match rate, zero false denials, a median latency of 23.29~ms, and a 95th percentile latency of 108.58~ms.
    \item \textbf{Concurrency Stress Load:} Under a 50-worker thread pool simulating concurrent autonomous agent workflows, AEGIS delivers 23.3 requests per second with a mean latency of 21.1~ms and zero database lock errors.
\end{enumerate}

Finally, AEGIS delivers a hardware-inspired, dark-mode Security Operations Center (SOC) dashboard built with Streamlit and Plotly, rendering real-time causal provenance DAGs, cryptographic approval desks, and live adversarial sandboxes. AEGIS proves that provable, sub-25ms zero-trust mediation is computationally viable for enterprise-grade autonomous AI multi-agent ecosystems.

\vspace{0.8cm}
\noindent
\textbf{Keywords:} Autonomous AI Agents, Information Flow Control (IFC), Confused Deputy Problem, Indirect Prompt Injection, Provenance Tracking, Dual-Lattice Model, Human-in-the-Loop Security, Zero-Trust Architecture, Write-Ahead Logging (WAL).

\newpage
